\documentclass[11pt,a4paper]{article}
\usepackage[utf8]{inputenc}
\usepackage[T1]{fontenc}
\usepackage[french]{babel}
\usepackage{lmodern}
\usepackage{geometry}
\usepackage{booktabs}
\usepackage{longtable}
\usepackage{hyperref}
\usepackage{amsmath}
\geometry{margin=2.2cm}

\title{Projet Big Data\\Analyse tabulaire / series temporelles avec Python}
\author{Lorenzini Thomas \and Baptiste DIDIER \and Antoine DUARTE}
\date{Avril 2026}

\begin{document}
\maketitle
\tableofcontents
\newpage

\section{Description du projet}
Ce projet vise a exploiter des donnees industrielles massives afin d'en extraire des informations utiles a la prise de decision. Le jeu de donnees principal est \texttt{extrusion.csv} (serie temporelle), accompagne de \texttt{stat.csv} (statistiques de variables).

\section{Objectifs}
Les objectifs couverts sont:
\begin{itemize}
  \item exploration des donnees (structure, motifs, anomalies),
  \item pretraitement (valeurs manquantes, types, dedoublonnage),
  \item visualisation (lineaire, dispersion, barres, heatmap, KDE),
  \item interpretation des resultats et derivation d'informations.
\end{itemize}

\section{Jeu de donnees et chargement}
\subsection{Source du jeu de donnees}
Le jeu de donnees utilise provient de Kaggle:\\
\href{https://www.kaggle.com/datasets/podsyp/find-a-defect-in-the-production-extrusion-line}{https://www.kaggle.com/datasets/podsyp/find-a-defect-in-the-production-extrusion-line}

\subsection{Pourquoi ce jeu de donnees}
Ce jeu de donnees a ete retenu pour trois raisons principales:
\begin{itemize}
  \item il s'agit d'une vraie serie temporelle industrielle, donc pertinente pour un projet big data applique;
  \item le volume est important (grand nombre d'observations et nombreuses variables capteurs), ce qui permet de pratiquer des techniques de pretraitement realistes;
  \item le contexte de ligne d'extrusion est propice a la detection d'anomalies et a l'analyse des relations entre mesures (pression, debit, epaisseur), en lien direct avec les objectifs du projet.
\end{itemize}

\subsection{Nature des donnees}
Le fichier \texttt{extrusion.csv} contient des mesures temporelles industrielles (pression, debit, epaisseur, etc.).

\subsection{Chargement}
Le chargement a ete realise avec \texttt{pandas.read\_csv()} en selectionnant un sous-ensemble de variables pertinentes (forte variance + variables metier), pour conserver une analyse lisible et robuste.

Le sous-ensemble de variables est un compromis entre:
\begin{itemize}
  \item representativite du procede,
  \item lisibilite des analyses et visualisations,
  \item cout de calcul raisonnable sur machine etudiante.
\end{itemize}

\section{Pretraitement}
\subsection{Qualite initiale}
Resultats mesures sur l'echantillon de travail:
\begin{itemize}
  \item dimensions initiales: \textbf{226536 lignes x 15 colonnes},
  \item valeurs manquantes avant traitement: \textbf{1845},
  \item doublons avant traitement: \textbf{0}.
\end{itemize}

\subsection{Techniques appliquees}
\begin{itemize}
  \item Conversion de types: conversion numerique explicite des colonnes de mesures.
  \item Valeurs manquantes: imputation par \textbf{mediane}.
  \item Justification de l'imputation mediane: plus robuste que la moyenne face aux valeurs extremes (outliers) frequentes en donnees industrielles.
  \item Doublons: verification puis suppression via \texttt{drop\_duplicates()}.
  \item Feature engineering: creation de \texttt{pressure\_diff}, \texttt{pressure\_roll\_mean\_30} et \texttt{hour}.
\end{itemize}

Apres traitement:
\begin{itemize}
  \item valeurs manquantes: \textbf{0},
  \item doublons: \textbf{0},
  \item dimensions finales: \textbf{226536 lignes x 19 colonnes}.
\end{itemize}

\section{EDA et visualisations}
\subsection{Statistiques sommaires}
Les statistiques descriptives ont ete calculees avec \texttt{df.describe()} directement dans le notebook.

Ces statistiques permettent d'identifier rapidement:
\begin{itemize}
  \item les plages de variation,
  \item les ecarts de dispersion entre variables,
  \item les premieres incoherences potentielles avant modelisation.
\end{itemize}

\subsection{Graphiques realises}
\begin{itemize}
  \item lineaire (tendance temporelle),
  \item dispersion (debit vs pression),
  \item barres (pression moyenne par heure),
  \item heatmap de correlation,
  \item histogramme + KDE,
  \item boxplot,
  \item visualisation interactive Plotly (bonus).
\end{itemize}

Les figures sont produites et interpretees directement dans le notebook, sans generation de fichiers externes additionnels.

\section{Valeurs aberrantes et anomalies}
\subsection{Outliers (IQR)}
Le nombre de valeurs aberrantes detectees sur la pression par la methode IQR est de \textbf{29865}. Ce volume eleve peut refleter des regimes de production differents et non uniquement des erreurs.

\subsection{Anomalies statistiques (Z-score, bonus)}
Le nombre d'anomalies detectees avec la regle $|z| > 3$ sur la pression est faible, ce qui suggere surtout des changements de regime plutot que des erreurs massives.

\section{Manipulation des donnees}
\subsection{Grouping et aggregation}
Les heures de plus forte pression moyenne sont: 20h, 19h, 13h, 21h, 12h.

\subsection{Filtrage et tri}
Les 5\% plus fortes pressions ont ete extraites puis triees, avec des maxima observes autour de 595.

\subsection{Analyse temporelle (bonus)}
La moyenne mobile sur 30 pas de temps a ete appliquee pour lisser le bruit et faire ressortir les tendances globales.

\section{Interpretation et conclusions}
\begin{itemize}
  \item Le pipeline de pretraitement assure une base propre (0 valeurs manquantes restantes) pour l'analyse.
  \item Les correlations positives entre pression et debit suggerent des interactions fortes des variables de process.
  \item Les variations horaires indiquent des plages operationnelles plus chargees, exploitables pour le pilotage qualite.
  \item Le grand nombre d'outliers IQR invite a distinguer les regimes operationnels des erreurs de capteurs.
\end{itemize}

En synthese, le jeu de donnees choisi est pertinent pour illustrer un cycle complet de science des donnees: acquisition, nettoyage, exploration, extraction d'informations et interpretation metier. Les resultats obtenus montrent qu'une approche methodique permet deja de degager des indicateurs exploitables pour le suivi de production et la qualite.

\end{document}
